% titel: Blatt 10, A1: Konjugierte Säuren und Basen
% skript: Mortimer Kap. 17 (Säuren und Basen, Brønsted)

% Teil I: Werkzeuge, kein Lösungsweg.

Nach Brønsted ist eine Säure ein Protonendonator, eine Base ein
Protonenakzeptor. Säure und Base eines Paares unterscheiden sich um genau
ein Proton.

\begin{werkzeugkasten}[Konjugierte Paare]
\begin{itemize}
  \item \textbf{Konjugierte Base} einer Säure HS: entsteht durch Abgabe eines
        Protons, $\mathrm{HS \rightarrow S^- + H^+}$. Sie trägt eine
        Ladungseinheit weniger (bzw. eine um Eins verringerte Ladung).
  \item \textbf{Konjugierte Säure} einer Base B: entsteht durch Aufnahme
        eines Protons, $\mathrm{B + H^+ \rightarrow HB^+}$.
  \item Neutrales H$_2$O kann beides sein (amphot, je nach Partner).
  \item Zweiwertige/mehrprotonige Systeme (wie HCO$_3^-$) haben je Stufe ein
        eigenes Paar.
\end{itemize}
\textit{Skript:} Mortimer Kap. 17 (konjugierte Säure-Base-Paare).
\end{werkzeugkasten}

\begin{achtung}[Darauf achten]
\begin{itemize}
  \item Bei der konjugierten Base \emph{Ladung} zählen: HCO$_3^-$ gibt H$^+$
        ab $\Rightarrow$ CO$_3^{2-}$.
  \item Bei der konjugierten Säure ein Proton \emph{dazugeben}: SO$_4^{2-}$
        wird HSO$_4^-$, H$_2$O wird H$_3$O$^+$.
  \item Teilchen, die selbst ein H abgeben können (PH$_4^+$, H$_2$S) sind
        Säuren.
\end{itemize}
\end{achtung}
